% GENERATED FILE. Edit canonical lessons, then run npm run generate:books.
% canonical-source-hash: fnv1a64:ebf3104747f83137
% canonical-lessons: LA-C200-tollo, LA-C200-pervenio, LA-C200-agnosco, LA-C200-recuso, LA-C200-volito

\chapter{To lift, To reach, To recognise, To refuse, To fly about}
\label{ch:la-to-lift-to-reach-to-recognise-to-refuse-to-fly-about}
\hlchaptermodality{200}

\begin{chapteropening}
\textbf{By the end of this chapter:} \emph{I can say to lift, to reach, to recognise, to refuse and to fly about.}

Complete the last lesson of chapter 200: I can say to lift, to reach, to recognise, to refuse and to fly about.
\end{chapteropening}

\section[tollō, tollere]{tollō, tollere — to lift, to raise; to remove}

\label{lesson:LA-C200-tollo}



\begin{quote}
Before the new one: say the Latin for to share, then the Latin for to avoid.
\end{quote}

\subsection*{You'll want to know: tollō, tollere}
\textbf{tollō, tollere} — \textquotedblleft{}to lift, to raise; to remove\textquotedblright{}.

\begin{cousinweb}
\textbf{tollō} is inherited from Proto-Italic \emph{*tolnō}, \textquotedblleft{}to lift up\textquotedblright{}, usually traced to Proto-Indo-European \emph{*tl-neh2-}, \textquotedblleft{}to lift, to bear\textquotedblright{} (EtymDB 2.1); Greek \emph{talanton}, \textquotedblleft{}a balance, a weight\textquotedblright{}, is a cousin (Lewis \& Short).
\end{cousinweb}

\begin{grammarlens}[title={What you've built}]
One more word for this chapter.
\end{grammarlens}

\subsection*{Your turn}
\begin{itemize}
\raggedright
  \item \emph{Say it:} \emph{tollō, tollere}
  \item \emph{Say it:} \emph{tollō, tollere}, once more
  \item \emph{Say it:} say \emph{tollō, tollere}
  \item \emph{Recall:} say the Latin for to share, then the Latin for to avoid, then say \emph{tollō, tollere} again
\end{itemize}

\begin{tcolorbox}[breakable,colback=teal!4,colframe=teal!35!black,title={Before you move on}]
What does tollō, tollere mean? (\textquotedblleft{}To lift, to raise; to remove\textquotedblright{}.) Say it once more.
\end{tcolorbox}
\section[perveniō, pervenīre]{perveniō, pervenīre — to reach, to arrive at}

\label{lesson:LA-C200-pervenio}



\begin{quote}
Before the new one: say the Latin for to avoid, then the Latin for to lift.
\end{quote}

\subsection*{You'll want to know: perveniō, pervenīre}
\textbf{perveniō, pervenīre} — \textquotedblleft{}to reach, to arrive at\textquotedblright{}.

\begin{cousinweb}
\textbf{perveniō} is \emph{per}, \textquotedblleft{}through\textquotedblright{}, with \emph{veniō}, \textquotedblleft{}to come\textquotedblright{} (EtymDB 2.1): to come all the way.
\end{cousinweb}

\begin{grammarlens}[title={What you've built}]
One more word for this chapter.
\end{grammarlens}

\subsection*{Your turn}
\begin{itemize}
\raggedright
  \item \emph{Say it:} \emph{perveniō, pervenīre}
  \item \emph{Say it:} \emph{perveniō, pervenīre}, once more
  \item \emph{Say it:} say \emph{perveniō, pervenīre}
  \item \emph{Recall:} say the Latin for to avoid, then the Latin for to lift, then say \emph{perveniō, pervenīre} again
\end{itemize}

\begin{tcolorbox}[breakable,colback=teal!4,colframe=teal!35!black,title={Before you move on}]
What does perveniō, pervenīre mean? (\textquotedblleft{}To reach, to arrive at\textquotedblright{}.) Say it once more.
\end{tcolorbox}
\section[agnōscō, agnōscere]{agnōscō, agnōscere — to recognise}

\label{lesson:LA-C200-agnosco}



\begin{quote}
Before the new one: say the Latin for to lift, then the Latin for to reach.
\end{quote}

\subsection*{You'll want to know: agnōscō, agnōscere}
\textbf{agnōscō, agnōscere} — \textquotedblleft{}to recognise\textquotedblright{}.

\begin{cousinweb}
\textbf{agnōscō} is \emph{ad}, \textquotedblleft{}to\textquotedblright{}, with \emph{gnōscō}, \textquotedblleft{}to get to know\textquotedblright{}, the older form of nōscō (Lewis \& Short): the g of the old form survives here.
\end{cousinweb}

\begin{grammarlens}[title={What you've built}]
One more word for this chapter.
\end{grammarlens}

\subsection*{Your turn}
\begin{itemize}
\raggedright
  \item \emph{Say it:} \emph{agnōscō, agnōscere}
  \item \emph{Say it:} \emph{agnōscō, agnōscere}, once more
  \item \emph{Say it:} say \emph{agnōscō, agnōscere}
  \item \emph{Recall:} say the Latin for to lift, then the Latin for to reach, then say \emph{agnōscō, agnōscere} again
\end{itemize}

\begin{tcolorbox}[breakable,colback=teal!4,colframe=teal!35!black,title={Before you move on}]
What does agnōscō, agnōscere mean? (\textquotedblleft{}To recognise\textquotedblright{}.) Say it once more.
\end{tcolorbox}
\section[recūsō, recūsāre]{recūsō, recūsāre — to refuse}

\label{lesson:LA-C200-recuso}



\begin{quote}
Before the new one: say the Latin for to reach, then the Latin for to recognise.
\end{quote}

\subsection*{You'll want to know: recūsō, recūsāre}
\textbf{recūsō, recūsāre} — \textquotedblleft{}to refuse\textquotedblright{}.

\begin{cousinweb}
\textbf{recūsō} is built on \emph{causa}, \textquotedblleft{}a case, a reason\textquotedblright{} (Lewis \& Short), with \emph{re-}, \textquotedblleft{}back\textquotedblright{}: to plead against, to refuse.
\end{cousinweb}

\begin{grammarlens}[title={What you've built}]
One more word for this chapter.
\end{grammarlens}

\subsection*{Your turn}
\begin{itemize}
\raggedright
  \item \emph{Say it:} \emph{recūsō, recūsāre}
  \item \emph{Say it:} \emph{recūsō, recūsāre}, once more
  \item \emph{Say it:} say \emph{recūsō, recūsāre}
  \item \emph{Recall:} say the Latin for to reach, then the Latin for to recognise, then say \emph{recūsō, recūsāre} again
\end{itemize}

\begin{tcolorbox}[breakable,colback=teal!4,colframe=teal!35!black,title={Before you move on}]
What does recūsō, recūsāre mean? (\textquotedblleft{}To refuse\textquotedblright{}.) Say it once more.
\end{tcolorbox}
\section[volitō, volitāre]{volitō, volitāre — to fly about}

\label{lesson:LA-C200-volito}



\begin{quote}
Before the new one: say the Latin for to recognise, then the Latin for to refuse.
\end{quote}

\subsection*{You'll want to know: volitō, volitāre}
\textbf{volitō, volitāre} — \textquotedblleft{}to fly about\textquotedblright{}.

\begin{cousinweb}
\textbf{volitō} is a frequentative of \emph{volō}, \textquotedblleft{}to fly\textquotedblright{} (Lewis \& Short): to keep flying, to flit about. It is a different verb from \emph{volō}, \textquotedblleft{}to want\textquotedblright{}.
\end{cousinweb}

\begin{grammarlens}[title={What you've built}]
One more word for this chapter.
\end{grammarlens}

\subsection*{Your turn}
\begin{itemize}
\raggedright
  \item \emph{Say it:} \emph{volitō, volitāre}
  \item \emph{Say it:} \emph{volitō, volitāre}, once more
  \item \emph{Say it:} say \emph{volitō, volitāre}
  \item \emph{Recall:} say the Latin for to recognise, then the Latin for to refuse, then say \emph{volitō, volitāre} again
\end{itemize}

\begin{tcolorbox}[breakable,colback=teal!4,colframe=teal!35!black,title={Before you move on}]
What does volitō, volitāre mean? (\textquotedblleft{}To fly about\textquotedblright{}.) Say it once more.
\end{tcolorbox}
